% ==============================================================================
% SECCIÓN 10: CONTROL DE VERSIONES
% ==============================================================================
\section{Control de Versiones (Método y Herramientas)}

Utilizaremos la herramienta \textbf{GitHub} para el control de versiones, actualizando el repositorio compartido en cada sesión de prácticas, llevando a cabo un control supervisado y documentado correctamente, especificando qué cambios se han llevado a cabo.

\subsection{Estrategia y Disciplina en el Repositorio}
\begin{itemize}[leftmargin=*]
  \item \textbf{Actualización Sistemática:} Es obligatorio realizar commit y push de los avances generados al término de cada sesión presencial de prácticas, documentando de forma el trabajo efectuado.
  \item \textbf{Control Supervisado y Pull Requests:} Las ramas funcionales de trabajo deberán ser integradas mediante Pull Requests revisadas antes de su incorporación a la rama principal de desarrollo para evitar conflictos.
  \item \textbf{Prohibición de Código Roto:} Ningún cambio se integrará en la rama principal sin haber sido revisado por el equipo, especialmente por el \textbf{Gestor de Calidad}. Cuando existan pruebas unitarias, deberá comprobarse mediante \texttt{flutter test} que todas se ejecutan satisfactoriamente sobre el código modificado.
\end{itemize}

\subsection{Convención de Commits (Conventional Commits)}
Formato obligatorio: \texttt{<tipo>: <descripción concisa>}
\begin{itemize}[leftmargin=*]
  \item \texttt{feat}: Nueva funcionalidad implementada.
  \item \texttt{fix}: Corrección de un defecto o fallo.
  \item \texttt{docs}: Actualización exclusiva de documentación o actas.
  \item \texttt{style}: Cambios estéticos o de formateo sin impacto lógico.
  \item \texttt{refactor}: Modificación de código sin alteración funcional.
  \item \texttt{test}: Incorporación o modificación de pruebas unitarias/integración.
  \item \texttt{a11y}: Modificaciones específicas de accesibilidad (contrastes, VoiceOver/TalkBack).
  \item \texttt{chore}: Mantenimiento de configuración, dependencias o CI/CD.
\end{itemize}
